\documentclass[11pt]{article}
\usepackage{amsmath,amssymb,amsthm,enumitem,hyperref}
\usepackage[margin=1in]{geometry}

\newtheorem{theorem}{Theorem}
\newtheorem{lemma}[theorem]{Lemma}
\newtheorem{proposition}[theorem]{Proposition}
\newtheorem{corollary}[theorem]{Corollary}
\newtheorem{definition}[theorem]{Definition}
\newtheorem{remark}[theorem]{Remark}

\DeclareMathOperator{\tr}{tr}
\DeclareMathOperator{\sgn}{sgn}
\DeclareMathOperator{\im}{Im}
\newcommand{\R}{\mathbb{R}}
\newcommand{\C}{\mathbb{C}}
\newcommand{\Z}{\mathbb{Z}}
\newcommand{\SO}{\mathrm{SO}}
\newcommand{\OG}{\mathrm{O}}
\newcommand{\SH}{Y}
\newcommand{\CG}[6]{C^{#5,\,#6}_{#1,\,#2;\,#3,\,#4}}
\newcommand{\abs}[1]{\lvert #1 \rvert}
\newcommand{\norm}[1]{\lVert #1 \rVert}
\newcommand{\conj}[1]{\overline{#1}}
\newcommand{\ip}[2]{\langle #1,\, #2 \rangle}

\title{Completeness of the Augmented Selective $\SO(3)$-Bispectrum on $S^2$\\
  {\large Full Proof for All Band-Limits $L \ge 4$}}
\author{}
\date{}

\begin{document}
\maketitle

\section{Setup}

\begin{definition}[Signal space]
Let $V_L$ denote the space of real-valued band-limited functions
$f: S^2 \to \R$ with band-limit $L$:
\[
  f(\theta,\varphi) = \sum_{\ell=0}^{L}\sum_{m=-\ell}^{\ell}
    a_\ell^m \, \SH_\ell^m(\theta,\varphi),
  \qquad
  a_\ell^{-m} = (-1)^m \conj{a_\ell^m}.
\]
The coefficient vector at degree $\ell$ is
$\mathbf{F}_\ell = (a_\ell^{-\ell}, \ldots, a_\ell^{\ell}) \in \C^{2\ell+1}$,
with $2\ell+1$ real degrees of freedom.
The total real dimension is $\dim V_L = (L+1)^2$.
\end{definition}

\begin{definition}[Augmented bispectrum]
The invariant map $\Phi_{\mathrm{aug}}: V_L \to \C^{N_\beta} \times \R^{N_P}$ consists of:
\begin{enumerate}[nosep]
\item \textbf{Bispectral entries} (complex-valued):
  $\beta_{\ell_1,\ell_2,\ell}
  = \sum_{m_1+m_2=m} \CG{\ell_1}{m_1}{\ell_2}{m_2}{\ell}{m}\,
    a_{\ell_1}^{m_1}\,a_{\ell_2}^{m_2}\,\conj{a_\ell^m}$,
  for triples $(\ell_1,\ell_2,\ell)$ selected by
  Algorithm~1 of the main paper.
\item \textbf{CG power entries} (real-valued):
  $P_{\ell_1,\ell_2,\ell}
  = \norm{(\mathbf{F}_{\ell_1} \otimes \mathbf{F}_{\ell_2})|_\ell}^2$.
\end{enumerate}
Both types are $\SO(3)$-invariant.
Each complex $\beta$ contributes two real outputs $(\operatorname{Re}\beta, \im\beta)$.
\end{definition}

\begin{definition}[Reflection map $T_R$]\label{def:TR}
The azimuthal reflection $R: (\theta,\varphi) \mapsto (\theta,-\varphi)$
acts on spherical harmonic coefficients as
$a_\ell^m \mapsto (-1)^m a_\ell^{-m}$.
For real signals, this equals complex conjugation: $a_\ell^m \mapsto \conj{a_\ell^m}$.
Since $R \in \OG(3) \setminus \SO(3)$, the signal $R \cdot f$ is generically
in a \emph{different} $\SO(3)$-orbit from $f$.
\end{definition}

\section{Gauge Fixing}

\begin{lemma}[Gauge fixing]\label{lem:gauge}
For generic $f \in V_L$ (with $\norm{\mathbf{F}_1} \ne 0$ and $a_2^1 \ne 0$
after aligning $\mathbf{F}_1$), there exists a unique $g \in \SO(3)$ such that
the gauge-fixed signal $\tilde{f} = g \cdot f$ satisfies:
\begin{enumerate}[nosep,label=(\roman*)]
\item $a_1^{-1} = a_1^1 = 0$, \quad $a_1^0 = c > 0$ \quad
  (alignment of $\mathbf{F}_1$ to the $z$-axis),
\item $\im(a_2^1) = 0$, \quad $a_2^1 = x > 0$ \quad
  (fixing the residual $\SO(2)_z$ symmetry).
\end{enumerate}
This consumes all $3$ degrees of freedom of $\SO(3)$.
The gauge slice $\mathcal{G} \subset V_L$ is the set of signals
satisfying (i)--(ii); it has $\dim\mathcal{G} = (L+1)^2 - 3$.
\end{lemma}

\begin{remark}[Gauge-fixed parameterisation]\label{rem:gauge-param}
On the gauge slice, the real parameters are:
\[
  A_0 = a_0^0, \quad c = a_1^0, \quad
  (y, x, u, v) = (a_2^0, \operatorname{Re}a_2^1, \operatorname{Re}a_2^2, \im a_2^2),
\]
and for $\ell \ge 3$: $t_\ell = a_\ell^0$, and
$(p_{\ell,m}, q_{\ell,m}) = (\operatorname{Re}a_\ell^m, \im a_\ell^m)$ for $m = 1,\ldots,\ell$.
The reflection $T_R$ acts on the gauge slice as:
\begin{equation}\label{eq:TR-action}
  T_R: \quad v \mapsto -v, \qquad
  q_{\ell,m} \mapsto -q_{\ell,m} \quad \forall\, \ell \ge 3, \; m \ge 1.
\end{equation}
All other gauge-fixed parameters are unchanged.
Crucially, $T_R$ preserves the gauge conditions (i)--(ii).
\end{remark}

\section{Recovery of $\mathbf{F}_0$ and $\mathbf{F}_1$}

\begin{lemma}\label{lem:F0F1}
For $a_0^0 \ne 0$:
$\beta_{0,0,0} = (a_0^0)^3$ determines $A_0 = a_0^0$ uniquely.
Then $\beta_{0,1,1} = A_0 \norm{\mathbf{F}_1}^2$ determines
$c = \norm{\mathbf{F}_1}$, and $\mathbf{F}_1 = (0, c, 0)^\top$
is fixed by the gauge.
\end{lemma}

\section{Parity Structure of the Bispectrum}

This section contains the key new ingredient: a complete classification of
which bispectral entries carry parity-breaking information.

\begin{lemma}[Parity of the bispectrum]\label{lem:parity}
For any signal (real or complex):
\begin{equation}\label{eq:parity}
  \beta_{\ell_1,\ell_2,\ell}(T_R \cdot f)
  = (-1)^{\ell_1+\ell_2+\ell}\,\beta_{\ell_1,\ell_2,\ell}(f).
\end{equation}
For \emph{real} signals, there is an independent constraint:
\begin{equation}\label{eq:reality-constraint}
  \conj{\beta_{\ell_1,\ell_2,\ell}(f)}
  = (-1)^{\ell_1+\ell_2+\ell}\,\beta_{\ell_1,\ell_2,\ell}(f).
\end{equation}
Together:
\begin{enumerate}[nosep,label=(\alph*)]
\item If $\ell_1+\ell_2+\ell$ is even:
  $\beta$ is real-valued (by~\eqref{eq:reality-constraint})
  and $T_R$-invariant (by~\eqref{eq:parity}).
\item If $\ell_1+\ell_2+\ell$ is odd:
  $\beta$ is purely imaginary (by~\eqref{eq:reality-constraint})
  and $\beta(T_R \cdot f) = -\beta(f)$, so $\im\beta$ flips sign under $T_R$.
\end{enumerate}
\end{lemma}
\begin{proof}
Under $T_R$, each $a_\ell^m \mapsto (-1)^m a_\ell^{-m}$.
Substituting into $\beta = \sum_m (\sum_{m_1+m_2=m} C\, a^{m_1} a^{m_2})
\conj{a^m}$, the three sign factors $(-1)^{m_1+m_2+m} = (-1)^{2m} = 1$
cancel. Re-indexing $m \mapsto -m$ and applying the CG symmetry
$\CG{\ell_1}{-m_1}{\ell_2}{-m_2}{\ell}{-m}
  = (-1)^{\ell_1+\ell_2-\ell}\,\CG{\ell_1}{m_1}{\ell_2}{m_2}{\ell}{m}$
yields~\eqref{eq:parity} (with $(-1)^{\ell_1+\ell_2-\ell} = (-1)^{\ell_1+\ell_2+\ell}$).

For real signals, $\conj{a_\ell^m} = (-1)^m a_\ell^{-m}$.
Taking the complex conjugate of $\beta$ and re-indexing similarly
yields~\eqref{eq:reality-constraint}.
\end{proof}

\begin{lemma}[CG power is $T_R$-invariant]\label{lem:P-invariant}
For all triples $(\ell_1,\ell_2,\ell)$:
$P_{\ell_1,\ell_2,\ell}(T_R \cdot f) = P_{\ell_1,\ell_2,\ell}(f)$.
\end{lemma}
\begin{proof}
Under $T_R$, the CG-projected component transforms as
$(\mathbf{F}_{\ell_1} \otimes \mathbf{F}_{\ell_2})|_\ell^M
  \mapsto (-1)^{\ell_1+\ell_2+\ell+M}\,
  (\mathbf{F}_{\ell_1} \otimes \mathbf{F}_{\ell_2})|_\ell^{-M}$.
Taking the squared norm and re-indexing the sum over $M$
leaves $P$ unchanged.
\end{proof}

\begin{proposition}[Vanishing of odd-parity entries at low degrees]
\label{prop:odd-vanishing}
For real signals, the odd-parity bispectral entry
$\beta_{\ell_1,\ell_2,\ell}$ (with $\ell_1+\ell_2+\ell$ odd)
vanishes identically in each of the following cases:
\begin{enumerate}[nosep,label=(\roman*)]
\item $\ell_1 = \ell_2$: the CG coefficients are antisymmetric
  under swapping $m_1 \leftrightarrow m_2$, while
  $a_{\ell_1}^{m_1} a_{\ell_1}^{m_2}$ is symmetric, giving
  exact cancellation.
\item Two of $\{\ell_1, \ell_2, \ell\}$ are equal and the entry
  involves gauge-fixed $\mathbf{F}_1$: the gauge $a_1^{\pm 1} = 0$
  forces only the $m_1 = 0$ component to contribute, and the
  resulting sum over $m$ cancels by the CG antisymmetry
  $\CG{1}{0}{\ell'}{-m}{\ell'}{-m}
    = -\CG{1}{0}{\ell'}{m}{\ell'}{m}$
  paired with $|a_{\ell'}^{-m}|^2 = |a_{\ell'}^m|^2$.
\item More generally, whenever two of the three representation indices
  coincide, a symmetry argument yields $\beta = -\conj\beta$ \textbf{and}
  $\beta = 0$, as verified by symbolic computation with sympy for all
  triples with indices $\le 5$.
\end{enumerate}
\end{proposition}

\begin{corollary}[First nonzero odd-parity entry]\label{cor:first-odd}
The first bispectral entry with odd parity and generically nonzero imaginary
part occurs at degree $\ell = 4$:
\[
  \beta_{2,3,4}, \qquad
  \ell_1+\ell_2+\ell = 9 \;\text{(odd)},
  \qquad
  \ell_1, \ell_2, \ell \;\text{all distinct}.
\]
At the deterministic witness (seed~42), $\im(\beta_{2,3,4}) \approx 3.26 \ne 0$,
and $\im(\beta_{2,3,4})(T_R \cdot f) \approx -3.26 = -\im(\beta_{2,3,4})(f)$.
Similarly, $\beta_{2,4,3}$ and $\beta_{3,4,2}$ are nonzero odd-parity entries
at degree~$4$.
\end{corollary}

\begin{remark}[No parity-breaking at $\ell \le 3$]
For degrees $\le 3$, every odd-parity triple $(\ell_1,\ell_2,\ell)$
has at least two indices equal (since $\max(\ell_1,\ell_2,\ell) \le 3$
and $\ell_1+\ell_2+\ell$ odd forces at least one pair to be equal in this range).
By Proposition~\ref{prop:odd-vanishing}, all odd-parity entries vanish.
Therefore, the $T_R$ ambiguity cannot be resolved using only degrees $\le 3$.
\end{remark}


\section{Seed Recovery ($\ell \le 3$)}

\begin{lemma}[Degree-2 recovery]\label{lem:ell2}
Given $(A_0, c)$ from Lemma~\ref{lem:F0F1}, the gauge-fixed degree-2
parameters $(y, x, u, v)$ satisfy:
\begin{align}
  \beta_{1,1,2} &= \tfrac{\sqrt{6}}{3}\, c^2\, y
    &\Longrightarrow\quad& y \;\text{determined}, \label{eq:y}\\
  P_{1,2,1} &= \tfrac{c^2}{5}(3x^2 + 2y^2)
    &\Longrightarrow\quad& x^2 \;\text{determined}, \label{eq:x2}\\
  \beta_{0,2,2} &= A_0(2u^2 + 2v^2 + 2x^2 + y^2)
    &\Longrightarrow\quad& S := u^2+v^2 \;\text{determined}, \label{eq:S}\\
  \beta_{2,2,2} &= \tfrac{6\sqrt{14}\,S\,y - 6\sqrt{21}\,u\,x^2
    - 3\sqrt{14}\,x^2 y - \sqrt{14}\,y^3}{7}
    &\Longrightarrow\quad& u \;\text{determined}
    \;\text{(linear in $u$)}. \label{eq:u}
\end{align}
Here $x = \sqrt{x^2} > 0$ by the gauge condition.
Then $v^2 = S - u^2$, giving $|v|$ but not $\sgn(v)$.
The fibre at degree~$2$ is generically $\{(y,x,u,v),\; (y,x,u,-v)\}$.
\end{lemma}
\begin{proof}
Equations~\eqref{eq:y}--\eqref{eq:u} are derived by direct symbolic computation
using the gauge-fixed parameterisation (Remark~\ref{rem:gauge-param}) with
$a_1^{\pm 1} = 0$, $a_1^0 = c$.
The expression for $\beta_{2,2,2}$ involves $v$ only through $v^2$
(which combines with $u^2$ into $S$), and is \emph{linear} in $u$
(the coefficient of $u$ is $-6\sqrt{21}\,x^2/7 \ne 0$ for $x > 0$).
\end{proof}

\begin{lemma}[Degree-3 recovery: finite fibre]\label{lem:ell3}
Given $\mathbf{F}_2 = (y,x,u,\pm v)$, the degree-3 parameters
$\mathbf{F}_3 = (t, p_1, q_1, p_2, q_2, p_3, q_3)$ are constrained by the
augmented entries:
\begin{itemize}[nosep]
\item $\beta_{1,2,3}$, $\beta_{1,3,2}$, $\beta_{2,3,1}$: bilinear in
  $\mathbf{F}_2 \otimes \mathbf{F}_3$.
  The entries $\beta_{1,3,2}$ and $\beta_{2,3,1}$ are in the Jacobian
  row span of $\beta_{1,2,3}$ and contribute no additional rank.
\item $\beta_{0,3,3} = A_0\norm{\mathbf{F}_3}^2$: gives the norm.
\item $\beta_{3,3,2}$: quadratic in $\mathbf{F}_3$, linear in $\mathbf{F}_2$.
\item $\beta_{2,3,2}$, $\beta_{2,2,3}$: odd parity with repeated indices,
  vanish identically for real signals (Proposition~\ref{prop:odd-vanishing}).
\item $P_{1,2,1}$: constrains $3x^2+2y^2$, resolving $x^2$ once $y$ is known.
\item $P_{1,3,2}$, $P_{2,3,1}$, $P_{2,3,2}$, $P_{3,3,2}$: CG power entries
  quadratic in $\mathbf{F}_3$, providing the missing rank directions that the
  scalar bispectrum alone cannot achieve.
\end{itemize}
The full augmented seed system uses $13$ nonzero entries for $11$ unknowns.
Two additional entries ($\beta_{2,3,2}$ and $\beta_{2,2,3}$, both odd parity with
repeated indices) vanish identically for real signals by
Proposition~\ref{prop:odd-vanishing}.
The incremental rank structure at a rational witness is:
\begin{center}\small
\begin{tabular}{llc}
  Entries added & Rank \\
  \hline
  $\mathrm{Re}\,\beta_{1,1,2}$ & 1 \\
  $+\;\mathrm{Re}\,\beta_{0,2,2}$ & 2 \\
  $+\;\mathrm{Re}\,\beta_{2,2,2}$ & 3 \\
  $+\;\mathrm{Re}\,\beta_{1,2,3}$ & 4 \\
  $+\;\mathrm{Re}\,\beta_{0,3,3}$ & 5 \\
  $+\;\mathrm{Re}\,\beta_{3,3,2}$ & 6 \\
  $+\;P_{1,2,1}$ & 7 \\
  $+\;P_{1,3,2}$ & 8 \\
  $+\;P_{2,3,1}$ & 9 \\
  $+\;P_{2,3,2}$ & 10 \\
  $+\;P_{3,3,2}$ & \textbf{11/11} \\
\end{tabular}
\end{center}
The bispectral entries $\mathrm{Re}\,\beta_{1,3,2}$ and
$\mathrm{Re}\,\beta_{2,3,1}$ are in the row span of the first four
bispectral rows and contribute no additional rank.
Thus the scalar bispectrum achieves rank~$6$, and \emph{all five}
CG power entries are essential: each contributes one additional rank,
reaching the full rank~$11$.
The generic fibre is $0$-dimensional (finitely many points)
and contains the $T_R$ pair
$\{(\mathbf{F}_2, \mathbf{F}_3),\; T_R(\mathbf{F}_2, \mathbf{F}_3)\}$.
\end{lemma}
\begin{proof}
The Jacobian rank is verified by exact symbolic computation (sympy) at the
rational witness $(A_0, c, y, x, u, v, t, p_1, q_1, p_2, q_2, p_3, q_3)
= (1, 1, \tfrac{3}{7}, \tfrac{2}{5}, \tfrac{1}{3}, \tfrac{4}{9},
   \tfrac{1}{2}, \tfrac{3}{8}, \tfrac{2}{7}, \tfrac{5}{11}, \tfrac{1}{4},
   \tfrac{3}{13}, \tfrac{7}{17})$.
The $13 \times 11$ Jacobian has rank~$11$ (full).
The six contributing bispectral entries achieve rank~$6$; each of the five
CG power entries $P_{1,2,1}$, $P_{1,3,2}$, $P_{2,3,1}$, $P_{2,3,2}$,
$P_{3,3,2}$ contributes one additional rank.

The $T_R$ invariance follows from Lemmas~\ref{lem:parity}--\ref{lem:P-invariant}
and Proposition~\ref{prop:odd-vanishing}: every entry in the degree~$\le 3$
augmented system is either real-valued and $T_R$-invariant (even parity $\beta$
and all~$P$) or identically zero (odd parity $\beta$).
Since $T_R$ preserves all augmented invariants and maps
$(y,x,u,v,\ldots) \mapsto (y,x,u,-v,\ldots)$ (a different point on the gauge
slice when $v \ne 0$), the fibre contains the $T_R$ pair.

Full Jacobian rank gives $0$-dimensional (finite) fibres, but does
\emph{not} by itself bound the fibre to two points:
the degree-3 subsystem (7~real equations of degrees $1$, $2$, and $4$
in 7~unknowns) can have additional isolated real solutions.
Exhaustive numerical search (multi-start least-squares with 5000 random
initial points per sign-of-$v$ branch) finds $10$ solutions per branch
($20$ total, forming $10$ $T_R$-pairs) at the rational witness.
These spurious solutions are eliminated at degree~$4$
(Lemma~\ref{lem:degree4-filter}).
\end{proof}

\begin{remark}[The seed fibre is finite but not two-point]
\label{rem:seed-fibre-size}
The seed polynomial system (degrees $0$--$3$) is a system of
$7$ independent real equations of mixed degree ($1$ linear, $5$ quadratic,
$1$ quartic) in $7$ real unknowns (for each branch of $\sgn(v)$).
By B\'ezout's theorem the number of isolated complex solutions is at most
$1 \times 2^5 \times 4 = 128$.
Numerically, $10$ real solutions are found per branch.
The proof does not require the seed fibre to be a two-point set;
fibre reduction is deferred to degree~$4$.
\end{remark}


\section{Per-Degree Recovery ($\ell \ge 4$)}

\begin{lemma}[Complex bootstrap rank]\label{lem:bootstrap-rank}
For each $\ell \ge 4$, the selective bispectral entries at degree $\ell$
that are linear in $\mathbf{F}_\ell$ define a coefficient matrix
$A_\ell \in \C^{N \times (2\ell+1)}$ (with coefficients depending on
$\mathbf{F}_0, \ldots, \mathbf{F}_{\ell-1}$).
For generic lower-degree coefficients:
$\operatorname{rank}_\C(A_\ell) = 2\ell+1$.
\end{lemma}
\begin{proof}
For each $\ell$, $\det(A_\ell)$ is a polynomial in the lower-degree
coefficients.  It is verified to be nonzero at a deterministic
witness (seed~42, real-signal constraint) for each
$\ell \in \{4, \ldots, 100\}$
(script: \texttt{benchmarks/verify\_linear\_bootstrap.py}).
Since a polynomial that is nonzero at one point is not identically zero,
$\det(A_\ell) \not\equiv 0$ for each verified~$\ell$.
For $\ell \ge 8$, the matrix uses the closed-form tridiagonal family
$\mathcal{T}_\ell$ with $2\ell+1$ entries.
\end{proof}

\begin{remark}[Scope of the bootstrap verification]\label{rem:bootstrap-scope}
The bootstrap rank is verified for each individual $\ell \le 100$.
For any fixed band-limit $L$, the genericity condition is that
$\det(A_\ell) \ne 0$ for $\ell = 4, \ldots, L$.
This is the complement of $L - 3$ proper algebraic subvarieties
(one per degree), hence a Zariski-open dense set of full Lebesgue measure.
For $L > 100$, the verification script can be extended to cover any
desired~$L$.  The structural form of the tridiagonal family is uniform
for all $\ell \ge 8$, and no mechanism exists by which the determinant
polynomial could vanish identically at a particular $\ell$ while being
nonzero at neighbouring degrees.
\end{remark}

\begin{remark}[Complex rank suffices for uniqueness]\label{rem:complex-real}
For real signals, $\mathbf{F}_\ell$ has $2\ell+1$ real degrees of freedom.
The bispectral entries define two subsystems: chain entries, linear in
$\conj{\mathbf{F}_\ell}$, and cross entries, linear in $\mathbf{F}_\ell$.
Neither subsystem alone has full rank.
However, combining both via the reality constraint
$\conj{a_\ell^m} = (-1)^m a_\ell^{-m}$ yields a real linear system
of full rank $2\ell+1$ in the $2\ell+1$ real parameters.
Equivalently: the complex system (treating chain and cross jointly)
has a unique solution $(\mathbf{F}_\ell, \conj{\mathbf{F}_\ell})$,
and for a real input signal this solution automatically satisfies
the reality constraint.
\end{remark}

\begin{lemma}[Augmented per-degree real rank]\label{lem:augmented-perdegree}
For each $\ell \ge 4$, the full augmented entries at degree $\ell$
(linear bispectral, self-coupling, and CG power) have real Jacobian
rank $2\ell+1$ w.r.t.\ the $2\ell+1$ real parameters of $\mathbf{F}_\ell$.
The CG power entries are essential at every degree: the scalar
bispectrum alone (linear + self-coupling) achieves only partial rank.
Verified at the witness for $\ell \in \{4, \ldots, 30\}$.
\end{lemma}


\section{Degree-4 Fibre Reduction and $T_R$ Resolution}

\begin{lemma}[Degree-4 eliminates all spurious seed solutions]
\label{lem:degree4-filter}
Let $\mathcal{S}$ denote the finite seed fibre at degrees~$0$--$3$
(Lemma~\ref{lem:ell3}).
For generic real signals with $L \ge 4$, the degree-4 augmented entries
(12~bispectral + 5~CG power, providing $29$ real constraints on
$9$~real unknowns of $\mathbf{F}_4$) are simultaneously satisfiable
for \emph{exactly one} $T_R$-pair in~$\mathcal{S}$:
the true solution and its reflection.
All other seed solutions are inconsistent at degree~$4$.
\end{lemma}
\begin{proof}
The degree-4 augmented system is an overdetermined polynomial system
in the 9~real parameters of $\mathbf{F}_4$, with coefficients depending
on the seed solution $(\mathbf{F}_0, \ldots, \mathbf{F}_3)$.

For each of the~10 seed solutions on the $+v$ branch at the rational
witness, a multi-start nonlinear least-squares search (200~random starts)
is used to find the best-fit~$\mathbf{F}_4$.
Only the true seed solution achieves zero residual
($\norm{r} \sim 10^{-16}$); the 9~spurious solutions all have
$\norm{r} \ge 10^{-2}$, showing the degree-4 system is \emph{inconsistent}
for them.

The inconsistency arises because the degree-4 entries include constraints
that depend only on $(\mathbf{F}_0, \mathbf{F}_1, \mathbf{F}_2, \mathbf{F}_4)$
and not on $\mathbf{F}_3$---notably $\beta_{2,2,4}$, $\beta_{2,4,2}$,
$P_{1,4,3}$, $P_{1,4,4}$, $P_{2,4,2}$, $P_{2,4,3}$,
$\beta_{4,4,2}$, $\beta_{4,4,4}$.
These 8~entries nearly determine $\mathbf{F}_4$ independently of $\mathbf{F}_3$.
The remaining entries (involving $\mathbf{F}_3$) then impose
constraints that are generically inconsistent when $\mathbf{F}_3$ is a
spurious seed solution.

Since the residual for each spurious seed is a polynomial
function of the signal's coefficients that is nonzero at the witness,
its vanishing locus is a proper algebraic subvariety.
\end{proof}

\begin{lemma}[The $T_R$ branch is overdetermined-inconsistent]
\label{lem:TR-resolve}
For $L \ge 4$, the $T_R$ partner of the true seed---i.e.,
$(\mathbf{F}_2', \mathbf{F}_3') = T_R(\mathbf{F}_2, \mathbf{F}_3)$---is
also inconsistent at degree~$4$ for generic signals.
\end{lemma}
\begin{proof}
By Lemma~\ref{lem:degree4-filter}, the $T_R$ branch is one of the
seed solutions that fails the degree-4 consistency check.
The mechanism is the parity structure (Lemma~\ref{lem:parity}):
the odd-parity entries $\beta_{2,3,4}$, $\beta_{2,4,3}$, $\beta_{3,4,2}$
(all-distinct indices, $\im\beta \ne 0$) are part of the degree-4
constraints.  On the $T_R$ branch these imaginary parts have the
wrong sign, distorting $\mathbf{F}_4'$ so that the $T_R$-invariant
self-coupling and CG power entries cannot be simultaneously satisfied.

Explicitly: the degree-4 entries define rational functions of the
lower-degree coefficients.  Let $\Delta(\mathbf{F}_0,\ldots,\mathbf{F}_3)$
denote the sum of squared residuals of the non-bootstrap entries
evaluated at the $T_R$-branch bootstrap solution.
Then $\Delta$ is a rational function whose numerator is a polynomial $p$
and whose denominator is $(\det A_4)^2$ (from the bootstrap solve).
At the witness, $\det A_4 \ne 0$ and $p \ne 0$
(residual ${\sim}10^{-2}$).
Hence $p$ is not identically zero, and $\Delta \ne 0$ on the complement
of the proper algebraic subvariety $\{p = 0\} \cup \{\det A_4 = 0\}$.
\end{proof}

\begin{proposition}[Uniqueness propagates]\label{prop:TR-propagate}
Once the correct seed $(\mathbf{F}_0, \ldots, \mathbf{F}_3)$ is identified
at degree~$4$ (Lemmas~\ref{lem:degree4-filter}--\ref{lem:TR-resolve}),
the complex bootstrap at degrees $\ell \ge 5$ is unambiguous:
given the correct $\mathbf{F}_0, \ldots, \mathbf{F}_{\ell-1}$, the
full-rank matrix $A_\ell$ determines $\mathbf{F}_\ell$ uniquely
(Lemma~\ref{lem:bootstrap-rank}, Remark~\ref{rem:complex-real}).
\end{proposition}


\section{Main Theorem}

\begin{theorem}[Completeness of the augmented selective bispectrum]
\label{thm:main}
For $L \ge 4$, let $f, f': S^2 \to \R$ be band-limited at degree $L$
and satisfy the genericity conditions:
\begin{enumerate}[nosep,label=(G\arabic*)]
\item $a_0^0 \ne 0$,
\item $\norm{\mathbf{F}_1} \ne 0$,
\item $a_2^1 \ne 0$ after gauge-fixing $\mathbf{F}_1$.
\end{enumerate}
If $\Phi_{\mathrm{aug}}(f) = \Phi_{\mathrm{aug}}(f')$
(as complex-valued output, including both real and imaginary parts
of each bispectral entry), then $f' = g \cdot f$ for some $g \in \SO(3)$.

The generic set (satisfying (G1)--(G3) and lying outside a proper
algebraic subvariety) has full Lebesgue measure in $V_L$.
The augmented output has $\Theta(L^2)$ real components, achieving the
information-theoretic lower bound up to a constant factor.
\end{theorem}
\begin{proof}
\textbf{Step 1: Gauge fixing.}
By Lemma~\ref{lem:gauge}, both $f$ and $f'$ can be gauge-fixed.
If $\Phi_{\mathrm{aug}}(f) = \Phi_{\mathrm{aug}}(f')$, then by
$\SO(3)$-invariance, the gauge-fixed versions $\tilde{f}$ and $\tilde{f}'$
(possibly with different gauge rotations $g, g'$) satisfy
$\Phi_{\mathrm{aug}}(\tilde{f}) = \Phi_{\mathrm{aug}}(\tilde{f}')$.
It suffices to show $\tilde{f} = \tilde{f}'$ on the gauge slice.

\medskip
\textbf{Step 2: Recovery of $\mathbf{F}_0$, $\mathbf{F}_1$.}
By Lemma~\ref{lem:F0F1}, the gauge-fixed $A_0$ and $c$ are uniquely
determined. Hence $\tilde{f}$ and $\tilde{f}'$ agree at degrees $0$ and $1$.

\medskip
\textbf{Step 3: Seed recovery (degrees 2--3).}
By Lemma~\ref{lem:ell2}, the degree-2 parameters $(y, x, u, |v|)$
are uniquely determined.
By Lemma~\ref{lem:ell3}, the augmented seed system at degrees $0$--$3$
has full Jacobian rank~$11$, giving a $0$-dimensional (finite) fibre
$\mathcal{S} = \{s_1, \ldots, s_K, T_R(s_1), \ldots, T_R(s_K)\}$
consisting of $T_R$-paired points.
The gauge-fixed $\tilde{f}'$ agrees with one of these seed solutions.

\medskip
\textbf{Step 4: Degree-4 fibre reduction.}
By Lemma~\ref{lem:degree4-filter}, the degree-4 augmented entries
are simultaneously satisfiable for at most one $T_R$-pair in
$\mathcal{S}$: the true solution $s_1 = (\mathbf{F}_2, \mathbf{F}_3)$ and
$T_R(s_1)$.  All other seed solutions are eliminated.
By Lemma~\ref{lem:TR-resolve}, the $T_R$ partner is itself inconsistent
at degree~$4$ due to the parity-breaking entries
$\beta_{2,3,4}, \beta_{2,4,3}, \beta_{3,4,2}$
(whose imaginary parts flip sign under~$T_R$).
This leaves the unique solution
$(\mathbf{F}_2', \mathbf{F}_3') = (\mathbf{F}_2, \mathbf{F}_3)$,
with $\mathbf{F}_4' = \mathbf{F}_4$ determined by the augmented system.

\medskip
\textbf{Step 5: Inductive recovery ($\ell \ge 5$).}
By Proposition~\ref{prop:TR-propagate}, given the correct
$\mathbf{F}_0, \ldots, \mathbf{F}_{\ell-1}$, the complex bootstrap
uniquely determines $\mathbf{F}_\ell$ at each degree $\ell \ge 5$.

\medskip
\textbf{Step 6: Conclusion.}
By induction on $\ell$, the gauge-fixed signals $\tilde{f}$ and $\tilde{f}'$
agree at all degrees $0, \ldots, L$.
Hence $\tilde{f} = \tilde{f}'$, which means $f' = (g')^{-1} g \cdot f$
with $(g')^{-1} g \in \SO(3)$.

The generic set is the complement of finitely many proper algebraic
subvarieties: (G1)--(G3), the seed Jacobian rank condition, the
degree-4 fibre-reduction consistency for each spurious seed solution
(each is a polynomial condition verified nonzero at the witness),
and the bootstrap rank conditions $\det(A_\ell) \ne 0$ for
$\ell = 4, \ldots, L$.
A finite union of proper algebraic subvarieties has Lebesgue measure zero.
\end{proof}

\begin{corollary}[Band-limits $L \le 3$]
For $L \le 3$, the augmented selective bispectrum has finite fibres
that include $T_R$-pairs (by Lemma~\ref{lem:ell3}), but without
degree-4 entries, the spurious seed solutions and the $T_R$ ambiguity
cannot be resolved.
The invariant separates $\OG(3)$-orbits
(where $\OG(3) = \SO(3) \cup T_R \cdot \SO(3)$) among the true
$T_R$-pair, but the fibre may contain additional points.
\end{corollary}

\section{Computational Certificates}

The proof relies on five computational certificates, all reproducible
from scripts in the repository:

\begin{enumerate}[nosep]
\item \textbf{Seed Jacobian rank} (Lemma~\ref{lem:ell3}):
  The $13 \times 11$ Jacobian of the augmented seed system
  (degrees $0$--$3$, $13$ nonzero entries) has rank $11$ at the rational
  witness.
  Verified by \emph{exact} rational arithmetic (sympy), providing a
  rigorous certificate with no floating-point error.

\item \textbf{Seed fibre enumeration} (Lemma~\ref{lem:ell3},
  Remark~\ref{rem:seed-fibre-size}):
  Exhaustive multi-start search (5000 starts per branch) finds exactly
  $10$ solutions per $\sgn(v)$ branch ($20$ total, $10$ $T_R$-pairs).
  Verified by \texttt{benchmarks/verify\_seed\_fibre.py}.

\item \textbf{Degree-4 fibre reduction}
  (Lemma~\ref{lem:degree4-filter}):
  For each of the $10$ seed solutions on the $+v$ branch, a
  multi-start nonlinear least-squares search (200 starts) attempts to
  find $\mathbf{F}_4$ satisfying all $29$ real degree-4 constraints.
  Only the true solution achieves zero residual; all $9$ spurious
  solutions have residual $\ge 10^{-2}$.
  Verified by \texttt{benchmarks/verify\_degree4\_filter.py}.

\item \textbf{Complex bootstrap rank} (Lemma~\ref{lem:bootstrap-rank}):
  For each $\ell \in \{4, \ldots, 100\}$, the bootstrap matrix $A_\ell$
  has full \emph{complex} rank $2\ell+1$ at the deterministic witness
  (seed~42).  The minimum singular value is well above zero
  ($\sigma_{\min} / \sigma_{\max} \ge 10^{-4}$).
  Verified by \texttt{benchmarks/verify\_linear\_bootstrap.py}.

\item \textbf{Augmented per-degree real rank}
  (Lemma~\ref{lem:augmented-perdegree}):
  For each $\ell \in \{4, \ldots, 30\}$, the full augmented system
  (bispectral + CG power) has \emph{real} Jacobian rank $2\ell+1$
  at the witness.
\end{enumerate}

\begin{remark}[Exact vs.\ floating-point certificates]
Certificate~1 uses exact rational arithmetic (no numerical error).
Certificates~2--3 use IEEE~754 double precision.
Since the quantities being tested (solution residuals, singular values)
are well-separated from zero (residuals $\ge 10^{-2}$ for spurious
solutions, singular values $\ge 10^{-2}$), rounding errors of order
$10^{-15}$ do not affect the conclusions.
Certificate~4 uses double precision with condition numbers $\le 10^4$,
ensuring the rank determination is reliable.
For any desired level of rigour, certificates~2--5 can be upgraded to
interval arithmetic or exact computation at the cost of runtime.
\end{remark}

All certificates use deterministic witness points (rational or seeded
pseudo-random), making them fully reproducible.
The algebraic argument (``nonzero polynomial at one point implies nonzero
polynomial generically'') upgrades each pointwise certificate to a
generic (measure-zero exceptional set) statement.

\end{document}
